\scanpage{44}

% B1 | text | top line and equality beneath right end
$\mathcal T(\bar{\mathbb F}_p)$ est un torseur canonique sous $H^1(\mathbb A,\mu_2)$

% B2 | display | top right
\[
\underset{\substack{\Vert\\\prod_\ell\mathbb Z_\ell^*/\mathbb Z_\ell^{*2}\\
\text{(y inclus $p$ et $\infty$)}}}{\vphantom{H^1(\mathbb A,\mu_2)}}
\]

% B3 | text | upper left, vertical note with brace
% margin note written sideways
\margin{\deleted{\unsure{à isom près}} On précisera l'objet : \unsure{question|quaterne|quotient} $(\unsure{\exists|\mathscr E}M,o)$ \uline{à isom près} (il y a deux classes d'isomorphie possibles si $(M,o)$ est l'une d'eux, l'autre est $(M^{(p)},o')$ que l'on en déduit en appliquant frobenius)}

% B4 | text | main upper paragraph
Ce torseur à priori dépend encore de $\bar{\mathbb F}_p$, \deleted{\unsure{mais il faut}} et cette
dépendance est précisée par la connaissance de l'opération
de $\hat{\mathbb Z}\ni\unsure{f|F}$ sur $\mathcal T(\bar{\mathbb F}_p)$. Cette opération se fait via
les opérations de $\hat{\mathbb Z}$ sur $C(A',T')$, \deleted{\unsure{$E$ étant tel}} \inserted{(\unsure{opér.|cycles} habituelles sur $A'(1)$, $f$ opère par $p$) et} $C(A_p,T_p)$
(via opération de $\deleted{\unsure{\hat{\mathbb Z}}}f$ sur $C(A_p,T_p)$ par frobenius) et
l'opération triviale sur $C(A_\infty,T_\infty)$. Bon on trouve que
$f\in\hat{\mathbb Z}$ opère par le $=$ élément $(p,f,p)$ que plus
haut, donc

% B5 | display | central isolated line
\[
f\text{ opère trivialement sur }\mathcal T(\bar{\mathbb F}_p)
\]

% B6 | text | central conclusion
Ainsi, on a un torseur canonique \inserted{$\mathcal T$} sous $H^1(\mathbb A,\mu_2)$.

% B7 | text | boxed mark and NB in lower half
$\square$ \uline{NB} Ceci suggère qu'on a une \deleted{section canonique} \inserted{gerbe triviale (\unsure{ou|au}}
\unsure{trivialisée|trivialisable|préhistorie}) de \deleted{\unsure{faisceaux}} $\mu_2$ sur $\mathbb A$, dont $\mathcal T$ serait
le torseur des classes d'isomorphie. On peut
en effet considérer les couples $(C,\deleted{\unsure{f}}\alpha)$, où
$C$ est une catégorie de Deuring de type $\unsure{\xi|\zeta}\in\mathrm{Br}(\mathbb Q)$, \inserted{d'où $T_C=T$,} et
$\alpha$ une \deleted{\unsure{équiv}} classe d'équivalence de classes d'isomorphie
de \inserted{l'équival.} $(C,T)_{\mathbb Z_p}$ avec $(C(\mathbb Z_p),T(\mathbb Z_p))$, modulo équivalences
provenant de $\underline{\underline{\mathrm{Pic}}}(\mathbb Z_p)$ ; i.e. une section du torseur de
groupe $\mathbb Z/2\mathbb Z$ déduit de celui de groupe $\underline{\mathrm{Equ}}(C(\mathbb Z_p),T(\mathbb Z_p))=\mathbb Z_p^*/\mathbb Z_p^{*2}\times\mathbb Z/2\mathbb Z$
déduit par extension du gpe structural à l'aide de $\mathrm{pr}_2$.

% B8 | text | lower left, sloping note connected to side note
% slanted
\margin{\unsure{trop} compliqué il suffit de \unsure{spécifier une|fichues} classe d'isom d'\unsure{isomorph.} d'équivalences de cat. de Deuring :
$C_{\mathbb Z_p}\simeq C(\mathbb Z_p)$}

% B9 | text | lower left, second note and cancellation
% slanted
\margin{en termes concrets : un $\alpha$ est donné (si $C$ défini par $o$) en \unsure{se} donnant une struct. de $o\otimes_{\mathbb Z}\mathbb Z_p$ sur l'anneau des endos $M_0$ ; d'un cristal sur $\bar{\mathbb F}_p$ (de rang $2$,
\deleted{pente $\frac12$ \unsure{et} \unsure{$\bigwedge^2$} \unsure{$A$} / \unsure{un choix d'un lien de Tate}}}

% B10 | text | bottom edge
Alors tout $C$ admet un $\alpha$, et deux $(C,\alpha),(C',\alpha')$
% archive sheet number: 43
